Expected initial norm matching gives 0/5 versus 5/5 generalization events at widths 64 and 256, respectively; mean final accuracies are 0.716 and 0.997.

Replacing ReLU with a square activation gives 5/5 versus 5/5 generalization events at widths 64 and 256, respectively; mean final accuracies are 0.986 and 0.984.

Increasing the prime to 47 gives 0/5 versus 5/5 generalization events at widths 64 and 256, respectively; mean final accuracies are 0.866 and 1.000.

Increasing the training fraction to 70\% gives 5/5 versus 5/5 generalization events at widths 64 and 256, respectively; mean final accuracies are 0.952 and 1.000.

Changing addition to multiplication gives 0/5 versus 5/5 generalization events at widths 64 and 256, respectively; mean final accuracies are 0.842 and 0.994.

Replacing arithmetic labels with independent random labels gives 0/5 versus 0/5 generalization events at widths 64 and 256, respectively; mean final accuracies are 0.033 and 0.035.

Changing the modulus fixes the training fraction, not the number of examples: prime 47 uses 1,104 training pairs rather than 480, and changes both the output dimension and the parameter count at fixed width. Cross-modulus waiting times are therefore not an isolated measure of arithmetic difficulty.

These are scope and intervention tests, with uncertainty reported in Table~\ref{tab:secondary}. In particular, norm matching changes the norm trend while retaining the width intervention, but also changes the initial logit scale. Random labels test a distinct target distribution; failure to generalize there is expected and is not additional evidence for a specific arithmetic circuit. Multiplication includes zero operands and does not have the same additive translation symmetry as addition.

At the reference setting, mean final addition-symmetry energy is 0.337 at width 64 and 0.622 at width 256 (Figure~\ref{fig:diagnostics}). This diagnostic is consistent with a change in output structure when it accompanies better accuracy, but its definition already uses the task symmetry and it cannot independently establish why a feature formed.

The exploratory direction-freezing intervention after 1,000 steps gives 0/5 versus 0/5 generalization events at widths 64 and 256. Mean final training accuracies are 1.000 and 1.000, and final held-out accuracies are 0.007 and 0.022, respectively. The width-256 trajectories are shown in Appendix Figure~\ref{fig:freezing}.

The exploratory direction-freezing intervention after 10,000 steps gives 0/5 versus 3/5 generalization events at widths 64 and 256. Mean final training accuracies are 1.000 and 1.000, and final held-out accuracies are 0.523 and 0.960, respectively. The width-256 trajectories are shown in Appendix Figure~\ref{fig:freezing}.

A step advantage need not be an advantage at a common compute proxy. For a budget of $64\times30{,}000$ hidden-unit updates, the actually evaluated checkpoints are width 64 at 30,000 steps, width 128 at 15,000 steps, width 256 at 7,500 steps. Their mean held-out accuracies are 0.716, 0.928, 0.830, respectively. This proxy reflects the width-dependent dense-layer arithmetic at fixed data size; it excludes fixed framework and optimizer overhead and is not measured wall-clock equality. Full intervals and per-seed values are supplied in the appendix.
